% Generated from the audited Markdown source; responsible author: Carptopus.
\documentclass[11pt]{article}
\usepackage[a4paper,margin=27mm]{geometry}
\usepackage[T1]{fontenc}
\usepackage[utf8]{inputenc}
\usepackage{lmodern}
\usepackage{microtype}
\usepackage{amsmath,amssymb,amsthm,mathtools}
\usepackage{needspace}
\usepackage{xcolor}
\usepackage[colorlinks=true,linkcolor=blue!55!black,citecolor=blue!55!black,urlcolor=blue!65!black]{hyperref}
\usepackage{fancyhdr}

\newtheorem{theorem}{Theorem}[section]
\newtheorem{lemma}[theorem]{Lemma}
\numberwithin{equation}{section}

\pagestyle{fancy}
\fancyhf{}
\fancyhead[L]{Carptopus}
\fancyhead[R]{Antichain-box polynomials}
\fancyfoot[C]{\thepage}
\setlength{\headheight}{14pt}
\setlength{\parindent}{0pt}
\setlength{\parskip}{0.45em}
\setlength{\emergencystretch}{4em}
\urlstyle{same}

\title{Real-rootedness, palindromicity, and $\gamma$-positivity\\
of antichain polynomials for
\texorpdfstring{$[2]\times[m]\times[n]$}{[2] x [m] x [n]}}
\author{Carptopus\\\href{mailto:carptopus@163.com}{carptopus@163.com}}
\date{26 August 2026}

\hypersetup{
  pdftitle={Real-rootedness, palindromicity, and gamma-positivity of antichain polynomials for [2] x [m] x [n]},
  pdfauthor={Carptopus},
  pdfsubject={Real-rootedness, palindromicity classification, and strict gamma-positivity for antichain generating polynomials of three-dimensional box posets},
  pdfkeywords={antichain generating polynomial, real-rootedness, palindromicity, gamma-positivity, poset, Jacobi polynomial, lattice path, tail switching, Zhang-Zhang polynomial}
}

\begin{document}
\maketitle

\begin{center}
\fcolorbox{black!35}{black!4}{\parbox{0.88\textwidth}{\centering\small
Internal candidate manuscript, version 0.2. The mathematical proof and complete
integrated manuscript have passed project-internal zero-trust audits. External
mathematical review and formal peer review are pending.}}
\end{center}

\begin{abstract}


Ding and Dong conjectured that the antichain generating polynomial of the product
poset $[2]\times[m]\times[n]$ is real-rooted for all positive integers $m,n$,
and that the adjacent-parameter family $[2]\times[r]\times[r+1]$ is
palindromic and $\gamma$-positive.
Independently, He, Langner, and Witek found by numerical experimentation a
$2\times2$ determinantal formula for the corresponding Zhang--Zhang polynomial,
which is obtained from the antichain polynomial by a unit variable shift, but
left the formula without a proof. We prove the translated determinantal formula
by a weight-preserving tail-switching
bijection on pairs of two-rowed arrays. The determinant then factors, after a
Möbius change of variables, into two controlled linear combinations of adjacent
Jacobi polynomials. Interlacing and endpoint signs show that every finite zero of
both factors lies in the required interval. Consequently the antichain generating
polynomial of $[2]\times[m]\times[n]$ has only negative real zeros for all positive
$m,n$, proving the first conjecture. We then prove that this polynomial is
palindromic if and only if $|m-n|=1$. In precisely these cases its
$\gamma$-coefficients are all strictly positive. This proves the
adjacent-parameter conjecture and verifies Ding and Dong's general
$\gamma$-positivity conjecture for the complete subfamily
$k=2$, $P=[m]\times[n]$.

\end{abstract}

\section{Introduction}


For a finite poset $P$, let

\nopagebreak[4]
\[
N_P(x)=\sum_{A\in\mathcal A(P)}x^{|A|}
\]

be its antichain generating polynomial, where the empty antichain is included.
Ding and Dong studied these polynomials for products involving minuscule posets
and proposed several positivity and real-rootedness conjectures \cite{ding2019}. One of them,
\cite[Conjecture~4.3]{ding2019}, states that

\nopagebreak[4]
\[
N_{[2]\times[m]\times[n]}(x)
\]

is real-rooted for every pair of positive integers $m,n$.

They further conjectured that

\nopagebreak[4]
\[
N_{[2]\times[r]\times[r+1]}(x)
\]

is palindromic and real-rooted for every $r\geq1$, and hence
$\gamma$-positive \cite[Conjecture~4.5]{ding2019}. Here a palindromic polynomial of degree
$D$ has a unique expansion

\nopagebreak[4]
\[
h(x)=\sum_{j=0}^{\lfloor D/2\rfloor}
\gamma_jx^j(1+x)^{D-2j}.
\tag{1.1}
\]

Following the strict convention in \cite{ding2019}, $\gamma$-positive means
$\gamma_j>0$ for every $j$. The $\gamma$-positivity conclusion of Conjecture
4.5 is the instance of their Conjecture B with $k=2$ and
$P=[r]\times[r+1]$; Conjecture 4.5 additionally asserts palindromicity and
real-rootedness.

The same family appears in the theory of extended strict order polynomials and
Zhang--Zhang polynomials of benzenoid systems. He, Langner, and Witek reported a
determinantal expression for the Zhang--Zhang polynomial of the hexagonal flake
$O(2,m,n)$ \cite{he2021}. Their formula was discovered through extensive numerical
experimentation and was explicitly stated without proof. The general relation
between extended strict order polynomials and the corresponding benzenoid
polynomials was developed in [3,4]. If $y$ is the Zhang--Zhang variable and $z$
is the antichain variable, their interface is $z=y+1$. We retain $x$ for the
antichain variable throughout the proof and make this translation explicit in
Section 2.

Our first result proves that determinant directly in the language of posets.
Define

\nopagebreak[4]
\[
A_{m,n}(x)=\sum_{j\geq0}\binom mj\binom njx^j,
\tag{1.2}
\]

\nopagebreak[4]
\[
B_{m,n}(x)=\sum_{r\geq0}
\binom{m+1}{r+2}\binom{n-1}{r}x^{r+1},
\tag{1.3}
\]

and

\nopagebreak[4]
\[
C_{m,n}(x)=\sum_{s\geq0}
\binom{m-1}{s}\binom{n+1}{s+2}x^{s+1}.
\tag{1.4}
\]

As usual, a binomial coefficient is zero when its lower index lies outside its
natural range.

\Needspace{0.30\textheight}
\begin{theorem}

For all positive integers $m,n$,

\nopagebreak[4]
\[
N_{[2]\times[m]\times[n]}(x)
=A_{m,n}(x)^2-B_{m,n}(x)C_{m,n}(x).
\tag{1.5}
\]

\end{theorem}

The proof is a signed tail-switching argument on two-rowed arrays. Its local move
is related to the crossing switch used by Krattenthaler \cite{krattenthaler1995}, but an independent
crosscut argument is needed to prove that every negative object has an admissible
inverse cut. This avoids the defective global crossing proof corrected by Rubey \cite{rubey}.

Our main result proves the conjecture and directly locates its zeros.

\Needspace{0.30\textheight}
\begin{theorem}

For all positive integers $m,n$, every zero of
$N_{[2]\times[m]\times[n]}(x)$ is real and negative.

\end{theorem}

The proof converts the two factors of (1.5) into linear combinations of adjacent
Jacobi polynomials. The coefficients fall exactly in the interval for which the
zeros remain in $[-1,1)$; the Möbius transformation back to $x$ sends the finite
zeros to $(-\infty,0)$.

\Needspace{0.30\textheight}
\begin{theorem}

For all positive integers $m,n$,

\nopagebreak[4]
\[
N_{[2]\times[m]\times[n]}(x)
\text{ is palindromic}
\quad\Longleftrightarrow\quad
|m-n|=1.
\]

For these and only these parameters it is strictly $\gamma$-positive.
Consequently \cite[Conjecture~4.5]{ding2019} holds, and \cite[Conjecture~B]{ding2019} holds for
$k=2$ and every rectangular minuscule poset $P=[m]\times[n]$.

\end{theorem}

The sufficiency is a parity identity for the Jacobi factorization. Necessity
also follows from the monicity classification \cite[Lemma~2.3(a)]{ding2019}; for
completeness we give a direct leading-coefficient proof. To obtain strict,
rather than merely nonnegative, $\gamma$-coefficients, we compute
$N_{[2]\times[r]\times[r+1]}(-1)$ explicitly and show that $-1$ is never a
zero.

The paper is organized as follows. Section 2 encodes antichains by two-rowed
arrays. Section 3 proves the tail-switching bijection and Theorem 1.1. Section 4
establishes the Jacobi factorization and proves Theorem 1.2. Section 5 proves
Theorem 1.3. Section 6 records the scope of the result and its relation to the
existing literature.

\section{Antichains and balanced two-rowed arrays}


Write $[r]=\{1,\ldots,r\}$ with its natural order. Within one fixed first-coordinate
layer of $[2]\times[m]\times[n]$, list the selected elements as

\nopagebreak[4]
\[
(a_1,c_1),\ldots,(a_p,c_p),
\]

where

\nopagebreak[4]
\[
a_1<\cdots<a_p,
\qquad
c_1>\cdots>c_p.
\]

Putting $b_i=n+1-c_i$ gives a balanced two-rowed array

\nopagebreak[4]
\[
U=\binom{a_1<\cdots<a_p}{b_1<\cdots<b_p},
\qquad
a_i\in[1,m],\quad b_i\in[1,n].
\tag{2.1}
\]

Conversely, every array of the form (2.1) uniquely determines an antichain in
$[m]\times[n]$. Hence the length-generating polynomial of one layer is
$A_{m,n}(x)$.

Let

\nopagebreak[4]
\[
U=\binom{a_1<\cdots<a_p}{b_1<\cdots<b_p},
\qquad
V=\binom{c_1<\cdots<c_q}{d_1<\cdots<d_q}
\tag{2.2}
\]

encode the two first-coordinate layers. An element from the first layer and an
element from the second layer are comparable precisely when some $i,j$ satisfy

\nopagebreak[4]
\[
a_i\leq c_j,
\qquad
b_i\geq d_j.
\tag{2.3}
\]

Shift the second array by setting

\nopagebreak[4]
\[
c'_j=c_j+1\in[2,m+1],
\qquad
d'_j=d_j-1\in[0,n-1].
\]

Then (2.3) is equivalent to the strict crossing condition

\nopagebreak[4]
\[
a_i<c'_j,
\qquad
d'_j<b_i.
\tag{2.4}
\]

Thus $A_{m,n}(x)^2$ counts all ordered pairs $(U,V)$, and the antichains of the
three-dimensional box are exactly the pairs having no crossing (2.4).

For later comparison, antichains also admit the following standard strict-map
interpretation. If $P$ is a finite poset, project an antichain in $[r]\times P$
onto an induced subposet $Q\subseteq P$. Its first coordinates give a strict
order-reversing map $Q\to[r]$, and this correspondence is bijective. Therefore

\nopagebreak[4]
\[
N_{[r]\times P}(x)=E^o_{P^*}(r,x),
\tag{2.5}
\]

where $E^o$ is the extended strict order polynomial. Since
$[2]\times[m]$ is self-dual, (2.5) identifies the polynomial considered here
with the relevant extended strict order polynomial. More precisely, using $z$
for the strict-order/antichain variable and $y$ for the Zhang--Zhang variable,
the equivalence theorem of Langner and Witek \cite{langner2022} gives

\nopagebreak[4]
\[
\operatorname{ZZ}(O(2,m,n),y)
=E^o_{[2]\times[m]}(n,y+1)
=N_{[2]\times[m]\times[n]}(y+1).
\tag{2.6}
\]

Thus Theorem 1.1 proves the determinant of \cite{he2021} after the explicit substitution
$x=y+1$; the two polynomials are not literally identical in the same variable.

\section{A two-gap tail-switching bijection}


We first describe the forward switch. Give a crossing $(i,j)$ of (2.4) the
coordinate

\nopagebreak[4]
\[
S(i,j)=(c'_j,b_i)
\]

and order these coordinates lexicographically. Let $(I,J)$ be the crossing with
maximal coordinate. Define

\nopagebreak[4]
\[
X_{\rm top}=(a_1,\ldots,a_I,c'_J,\ldots,c'_q),
\]

\nopagebreak[4]
\[
X_{\rm bot}=(b_1,\ldots,b_{I-1},d'_{J+1},\ldots,d'_q),
\tag{3.1}
\]

\nopagebreak[4]
\[
Y_{\rm top}=(c'_1,\ldots,c'_{J-1},a_{I+1},\ldots,a_p),
\]

and

\nopagebreak[4]
\[
Y_{\rm bot}=(d'_1,\ldots,d'_J,b_I,\ldots,b_p).
\tag{3.2}
\]

All four rows are strictly increasing. The two inequalities at the switch are
exactly (2.4). If $b_{I-1}\geq d'_{J+1}$, then $(I,J+1)$ would be a crossing with
larger first coordinate. If $c'_{J-1}\geq a_{I+1}$, then $(I+1,J)$ would be a
crossing with the same first coordinate and larger second coordinate. Empty
prefixes and suffixes handle the boundary cases.

The top row of $X$ has two more entries than its bottom row, while the bottom row
of $Y$ has two more entries than its top row. The corresponding generating
functions are $B_{m,n}(x)$ and $C_{m,n}(x)$. Indeed, if the shorter row of $X$
has length $r$, then its rows choose $r+2$ entries from $[1,m+1]$ and $r$ entries
from $[1,n-1]$ and carry weight $x^{r+1}$. The analogous statement for $Y$
gives (1.4). The switch preserves the union of the four rows and hence preserves
the weight

\nopagebreak[4]
\[
x^{\frac12(\text{sum of the four row lengths})}.
\]

It remains to prove that the switch is a bijection onto all pairs of arrays of
these two unbalanced types.

\Needspace{0.30\textheight}
\begin{lemma}[two-gap tail switch]

Let

\nopagebreak[4]
\[
X=\binom{\alpha_1<\cdots<\alpha_{r+2}}
          {\beta_1<\cdots<\beta_r},
\qquad
Y=\binom{\gamma_1<\cdots<\gamma_s}
          {\delta_1<\cdots<\delta_{s+2}},
\tag{3.3}
\]

where

\nopagebreak[4]
\[
\alpha\subseteq[1,m+1],\quad \beta\subseteq[1,n-1],\quad
\gamma\subseteq[2,m],\quad \delta\subseteq[0,n].
\]

Put

\nopagebreak[4]
\[
\gamma_0=1,\quad \gamma_{s+1}=m+1,
\qquad
\beta_0=0,\quad \beta_{r+1}=n.
\]

Then there is a pair $(i,j)\in[1,r+1]\times[1,s+1]$ such that

\nopagebreak[4]
\[
\alpha_i<\gamma_j,\qquad \gamma_{j-1}<\alpha_{i+1},
\tag{3.4}
\]

\nopagebreak[4]
\[
\beta_{i-1}<\delta_{j+1},\qquad \delta_j<\beta_i.
\tag{3.5}
\]

Call such a pair an admissible gap. Choose the admissible gap for which

\nopagebreak[4]
\[
T(i,j)=(\alpha_{i+1},\delta_{j+1})
\]

is lexicographically maximal. The reverse splice at this gap produces a unique
crossing pair $(U,V)$, and applying (3.1)--(3.2) at its maximal crossing returns
$(X,Y)$.

\end{lemma}

\begin{proof}

Set

\nopagebreak[4]
\[
\widehat\beta=(0,\beta_1,\ldots,\beta_r,n),
\qquad
\widehat\gamma=(1,\gamma_1,\ldots,\gamma_s,m+1).
\]

Inside the rectangle $R=[1,m+1]\times[0,n]$, join consecutive points of

\nopagebreak[4]
\[
P_i=(\alpha_i,\widehat\beta_i),\qquad 1\leq i\leq r+2,
\]

and of

\nopagebreak[4]
\[
Q_j=(\widehat\gamma_j,\delta_j),\qquad 1\leq j\leq s+2.
\]

Both coordinates increase strictly along every edge. The polyline $P$ joins the
bottom side of $R$ to its top side, while $Q$ joins the left side to the right
side. The rectangle crosscut lemma therefore implies that they intersect.

Choose incident edges at an intersection so that both their open horizontal
projections and their open vertical projections overlap. Such a choice is
immediate at an interior intersection. At a common vertex, take the two entering
edges or the two leaving edges; strict monotonicity gives the same conclusion.
If the polylines share a segment, use an interior point of that segment. At a
boundary intersection, the only apparently exceptional entering--leaving choice
would force one of the adjacent strictly increasing edges immediately outside
$R$, so it cannot occur. This also covers the four corners.

For the selected edge of $P$, say the edge from $P_i$ to $P_{i+1}$, and the
selected edge of $Q$, say the edge from $Q_j$ to $Q_{j+1}$, overlap of the two
open coordinate projections says

\nopagebreak[4]
\[
\alpha_i<\widehat\gamma_{j+1},\qquad
\widehat\gamma_j<\alpha_{i+1},
\]

\nopagebreak[4]
\[
\widehat\beta_i<\delta_{j+1},\qquad
\delta_j<\widehat\beta_{i+1}.
\]

Substitution of the augmented entries gives (3.4)--(3.5), proving that an
admissible gap exists.

At an admissible gap define

\nopagebreak[4]
\[
U_{\rm top}=(\alpha_1,\ldots,\alpha_i,
              \gamma_j,\ldots,\gamma_s),
\]

\nopagebreak[4]
\[
U_{\rm bot}=(\beta_1,\ldots,\beta_{i-1},
              \delta_{j+1},\ldots,\delta_{s+2}),
\tag{3.6}
\]

\nopagebreak[4]
\[
V'_{\rm top}=(\gamma_1,\ldots,\gamma_{j-1},
               \alpha_{i+1},\ldots,\alpha_{r+2}),
\]

and

\nopagebreak[4]
\[
V'_{\rm bot}=(\delta_1,\ldots,\delta_j,
               \beta_i,\ldots,\beta_r).
\tag{3.7}
\]

The four inequalities (3.4)--(3.5) are precisely the conditions that all rows
in (3.6)--(3.7) be strictly increasing. Undoing the shift,
$V_{\rm top}=V'_{\rm top}-1$ and $V_{\rm bot}=V'_{\rm bot}+1$, gives balanced
arrays in the ranges of (2.2). Moreover
$\alpha_i<\alpha_{i+1}$ and $\delta_j<\delta_{j+1}$ become a crossing of $U$
and $V'$ at the splice.

The local switch is the two-rowed-array operation (27) of Krattenthaler \cite{krattenthaler1995}
under the reindexing

\nopagebreak[4]
\[
M_1=V':\quad a_j=c'_j,\quad b_{j-1}=d'_j,
\qquad
M_2=U:\quad c_i=a_i,\quad d_i=b_i.
\]

Under this reindexing its crossing coordinate is both
$S=(c'_j,b_i)$ on the balanced side and
$T=(\alpha_{i+1},\delta_{j+1})$ on the unbalanced side. The four tails strictly
above and to the right of the splice are unchanged. Hence maximal $S$ is carried
to maximal $T$, and conversely; switching again at that coordinate restores all
four rows entry by entry. This proves the asserted inverse relation. Notice that
we use only the local switch and maximal-coordinate argument of \cite{krattenthaler1995}. The separate
crosscut argument above supplies existence on the unbalanced side, so the global
crossing proof corrected in \cite{rubey} is not used.

If $m=1$ or $n=1$, one of the interior rows may be empty, but each augmented
polyline still has at least one edge. The same crosscut and switch arguments
therefore include these boundary cases.
\end{proof}


\begin{proof}[Proof of Theorem 1.1]

The forward construction (3.1)--(3.2) and Lemma 3.1 give a weight-preserving
bijection between

\begin{enumerate}
\item the crossing pairs counted by $A_{m,n}(x)^2$, and
\item all unbalanced pairs counted by $B_{m,n}(x)C_{m,n}(x)$.

\end{enumerate}

Removing the crossing pairs from all pairs leaves exactly the antichains of
$[2]\times[m]\times[n]$. This proves (1.5).
\end{proof}


\section{Jacobi factorization and location of the zeros}


Because (1.5) is symmetric in $m,n$, assume $m\leq n$ and put

\nopagebreak[4]
\[
d=n-m,\qquad t=\frac{1+x}{1-x},\qquad R=m+n+1.
\tag{4.1}
\]

Let

\nopagebreak[4]
\[
G(x)={}_2F_1(1-m,1-n;3;x).
\]

Elementary binomial identities applied to (1.3)--(1.4) give

\nopagebreak[4]
\[
B_{m,n}(x)C_{m,n}(x)
=\binom{m+1}{2}\binom{n+1}{2}x^2G(x)^2.
\tag{4.2}
\]

The hypergeometric representation of Jacobi polynomials \cite[Eq.~18.5.8]{dlmf} yields

\nopagebreak[4]
\[
A_{m,n}(x)=(1-x)^mP_m^{(0,d)}(t),
\tag{4.3}
\]

and

\nopagebreak[4]
\[
G(x)=\frac{(1-x)^{m-1}}{\binom{m+1}{2}}
P_{m-1}^{(2,d)}(t).
\tag{4.4}
\]

Set

\nopagebreak[4]
\[
\rho=\sqrt{\frac{n(n+1)}{m(m+1)}}.
\]

Using $x/(1-x)=(t-1)/2$, equations (1.5) and (4.2)--(4.4) give

\nopagebreak[4]
\[
N_{[2]\times[m]\times[n]}(x)
=(1-x)^{2m}F_-(t)F_+(t),
\tag{4.5}
\]

where

\nopagebreak[4]
\[
F_\pm(t)=P_m^{(0,d)}(t)
\pm\frac{\rho}{2}(t-1)P_{m-1}^{(2,d)}(t).
\tag{4.6}
\]

Write $Q_k(t)=P_k^{(1,d)}(t)$. The adjacent-parameter identities for Jacobi
polynomials \cite[Eqs.~18.9.5--18.9.6]{dlmf}, after the symmetry
$t\mapsto-t$ and exchange of parameters, give

\nopagebreak[4]
\[
P_m^{(0,d)}(t)=\frac{n+1}{R}Q_m(t)-\frac nR Q_{m-1}(t),
\tag{4.7}
\]

and

\nopagebreak[4]
\[
(t-1)P_{m-1}^{(2,d)}(t)
=\frac{2m}{R}Q_m(t)-\frac{2(m+1)}R Q_{m-1}(t).
\tag{4.8}
\]

Consequently

\nopagebreak[4]
\[
RF_-(t)=\bigl(n+1-\rho m\bigr)Q_m(t)
+\bigl(\rho(m+1)-n\bigr)Q_{m-1}(t),
\tag{4.9}
\]

and

\nopagebreak[4]
\[
RF_+(t)=\bigl(n+1+\rho m\bigr)Q_m(t)
-\bigl(n+\rho(m+1)\bigr)Q_{m-1}(t).
\tag{4.10}
\]

We use the following elementary consequence of interlacing.

\Needspace{0.30\textheight}
\begin{lemma}

Let $d\geq0$ and $Q_k=P_k^{(1,d)}$. If

\nopagebreak[4]
\[
-\frac{m+1}{m}<c<\frac{m+d}{m},
\tag{4.11}
\]

then all $m$ zeros of $Q_m+cQ_{m-1}$ lie in $(-1,1)$. At either endpoint of
(4.11), the corresponding endpoint $1$ or $-1$ is a zero and all remaining
zeros lie in $(-1,1)$.

\end{lemma}

\begin{proof}

The Jacobi weight $(1-t)(1+t)^d$ is positive on $(-1,1)$, so the zeros of
$Q_m$ and $Q_{m-1}$ are simple and strictly interlace \cite[Sec.~18.2(vi)]{dlmf}.
Evaluating $Q_m+cQ_{m-1}$ at consecutive zeros of $Q_m$ gives $m-1$ zeros,
one in each intervening interval. The endpoint values are

\nopagebreak[4]
\[
Q_k(1)=k+1,
\qquad
Q_k(-1)=(-1)^k\binom{k+d}{k}.
\]

The endpoint values give the bounds

\nopagebreak[4]
\[
c>-\frac{m+1}{m}
\qquad\text{and}\qquad
c<\frac{m+d}{m}.
\]

If $c>0$, the $m-1$ interior sign changes are supplemented on the left: the
condition $c<(m+d)/m$ makes the signs at $-1$ and at the smallest zero of
$Q_m$ opposite. If $c<0$, the condition $c>-(m+1)/m$ similarly supplies the
last zero in the right outer interval. If $c=0$, the polynomial is $Q_m$
itself. Equality in the upper or lower bound makes $-1$ or $1$, respectively,
the final zero.
\end{proof}


We now apply the lemma. The coefficient ratio in (4.9) is

\nopagebreak[4]
\[
c_-=
\frac{\rho(m+1)-n}{n+1-\rho m}.
\tag{4.12}
\]

Both numerator and denominator are positive, as follows explicitly from

\nopagebreak[4]
\[
\rho^2(m+1)^2-n^2=\frac{n(m+n+1)}m>0
\]

and

\nopagebreak[4]
\[
(n+1)^2-\rho^2m^2
=\frac{(n+1)(m+n+1)}{m+1}>0.
\]

Moreover

\nopagebreak[4]
\[
c_-\leq\frac nm,
\]

because this inequality reduces to $\rho m\leq n$, whose square is equivalent
to $m\leq n$. Equality holds exactly when $m=n$. Thus all zeros of $F_-$ lie
in $(-1,1)$ when $m<n$; when $m=n$, one zero is $t=-1$ and the others remain
in $(-1,1)$.

The coefficient ratio in (4.10) is

\nopagebreak[4]
\[
c_+=-
\frac{n+\rho(m+1)}{n+1+\rho m}.
\tag{4.13}
\]

Clearly $c_+<n/m$, while

\nopagebreak[4]
\[
c_+>-\frac{m+1}{m}
\]

reduces, after cancellation of the common $\rho m(m+1)$ term, to
$mn<(m+1)(n+1)$. Hence every zero of $F_+$ lies in $(-1,1)$.

\begin{proof}[Proof of Theorem 1.2]

The inverse transformation in (4.1) is

\nopagebreak[4]
\[
x=\frac{t-1}{t+1},
\]

which maps $(-1,1)$ bijectively onto $(-\infty,0)$. Equations (4.5) and the
preceding analysis therefore place every finite zero of
$N_{[2]\times[m]\times[n]}(x)$ on the negative real axis. When $m=n$, the
endpoint zero $t=-1$ of $F_-$ corresponds to the point at infinity and accounts
for the degree drop from $2m$ to $2m-1$; it creates no finite zero. This proves
the theorem.
\end{proof}


\section{Palindromicity and strict \texorpdfstring{$\gamma$}{gamma}-positivity}


We now determine exactly when the polynomial in Theorem 1.2 is palindromic.
The sufficient direction comes from an additional symmetry of the Jacobi
factorization when the two rectangular parameters are adjacent.

\Needspace{0.30\textheight}
\begin{lemma}[adjacent-parameter parity]

For every positive integer $r$, put

\nopagebreak[4]
\[
J(t)=P_r^{(0,1)}(t),\qquad
\widetilde J(t)=P_r^{(1,0)}(t),
\]

\nopagebreak[4]
\[
H(t)=(t-1)P_{r-1}^{(2,1)}(t),\qquad
K(t)=(t+1)P_{r-1}^{(1,2)}(t).
\]

Then

\nopagebreak[4]
\[
J-\widetilde J=\frac12(H-K),
\tag{5.1}
\]

\nopagebreak[4]
\[
J+\widetilde J=\frac{r+2}{2r}(H+K).
\tag{5.2}
\]

Consequently

\nopagebreak[4]
\[
E_r(t)=J(t)^2-\frac{r+2}{4r}H(t)^2
\tag{5.3}
\]

is an even polynomial.

\end{lemma}

\begin{proof}

Set

\nopagebreak[4]
\[
u=\frac{t-1}{2},\qquad v=\frac{t+1}{2}.
\]

The Jacobi binomial expansion is

\nopagebreak[4]
\[
P_r^{(\alpha,\beta)}(t)
=\sum_{k=0}^r
\binom{r+\alpha}{k}
\binom{r+\beta}{r-k}u^{r-k}v^k.
\tag{5.4}
\]

With

\nopagebreak[4]
\[
T_k=\binom rk\binom{r+1}k,
\]

the coefficient of $u^{r-k}v^k$ in $J,\widetilde J,H,K$ is, respectively,

\nopagebreak[4]
\[
T_k\frac{r+1-k}{k+1},\qquad
T_k,\qquad
2T_k\frac{r-k}{k+1},\qquad
2T_k\frac{k}{k+1}.
\tag{5.5}
\]

Thus both sides of (5.1) have coefficient
$T_k(r-2k)/(k+1)$, and both sides of (5.2) have coefficient
$T_k(r+2)/(k+1)$. This includes $k=0,r$, where the appropriate coefficient
in $K$ or $H$ vanishes.

Jacobi symmetry gives

\nopagebreak[4]
\[
J(-t)=(-1)^r\widetilde J(t),\qquad
H(-t)=(-1)^rK(t).
\]

Therefore

\nopagebreak[4]
\[
\begin{aligned}
E_r(t)-E_r(-t)
&=(J-\widetilde J)(J+\widetilde J)\\
&\quad-\frac{r+2}{4r}(H-K)(H+K)=0
\end{aligned}
\]

by (5.1)--(5.2).
\end{proof}


\begin{proof}[Proof of Theorem 1.3]

The determinant (1.5) is symmetric in $m,n$, so assume $m\leq n$.
First let $n=r+1$ and $m=r$. Then $d=1$ and
$\rho^2=(r+2)/r$ in (4.5)--(4.6), giving

\nopagebreak[4]
\[
N_{[2]\times[r]\times[r+1]}(x)
=(1-x)^{2r}
E_r\!\left(\frac{1+x}{1-x}\right).
\tag{5.6}
\]

Under $x\mapsto1/x$, the argument of $E_r$ changes sign, while

\nopagebreak[4]
\[
\left(1-\frac1x\right)^{2r}=x^{-2r}(1-x)^{2r}.
\]

Lemma 5.1 therefore gives

\nopagebreak[4]
\[
x^{2r}N_{[2]\times[r]\times[r+1]}(1/x)
=N_{[2]\times[r]\times[r+1]}(x).
\tag{5.7}
\]

The coefficient of $x^{2r}$ in (1.5) is

\nopagebreak[4]
\[
(r+1)^2-r(r+2)=1,
\]

so the degree is exactly $2r$ and (5.7) proves palindromicity.

We next exclude every nonadjacent parameter pair. If $m=n$, the coefficient of
$x^{2m}$ cancels. The diagonal contribution to the coefficient of
$x^{2m-1}$ is $2m^2$, while the two off-diagonal contributions total
$2(m^2-1)$. Thus the degree is $2m-1$ and the leading coefficient is $2$.
Since the constant term is $1$, the polynomial is not palindromic.

If $n-m\geq2$, the degree is $2m$ and its leading coefficient is

\nopagebreak[4]
\[
L_{m,n}
=\binom nm^{\!2}
-\binom{n-1}{m-1}\binom{n+1}{m+1}
=\binom nm^{\!2}\frac{n-m}{n(m+1)}.
\tag{5.8}
\]

For $m=1$, this is $n(n-1)/2>1$. For $m\geq2$, write
$n=m+d$ with $d\geq2$. Then

\nopagebreak[4]
\[
\binom{m+d}{m}\geq\frac{(m+2)(m+1)}2,
\qquad
\frac{d}{m+d}\geq\frac2{m+2},
\]

and (5.8) yields

\nopagebreak[4]
\[
L_{m,n}\geq\frac{(m+2)(m+1)}2>1.
\]

Again the constant term is $1$, so the polynomial cannot be palindromic.
This completes the classification.

It remains to prove strict $\gamma$-positivity in the adjacent cases. Theorem
1.2 gives only negative real zeros. Because the polynomial is palindromic,
monic, has degree $2r$, and has constant term $1$, its zeros, counted with
multiplicity, can be paired as $-a_i,-a_i^{-1}$ with $a_i>0$. To show that the
resulting $\gamma$-coefficients are strictly positive, rather than merely
nonnegative, we must exclude $a_i=1$.

At $x=-1$, equation (4.3) gives

\nopagebreak[4]
\[
A_{r,r+1}(-1)=2^rJ(0).
\]

Coefficient extraction also gives

\nopagebreak[4]
\[
\begin{aligned}
A_{r,r+1}(-1)
&=\sum_j(-1)^j\binom rj\binom{r+1}j\\
&=[z^{r+1}](1-z)^r(1+z)^{r+1}\\
&=[z^{r+1}](1-z^2)^r(1+z).
\end{aligned}
\tag{5.9}
\]

Hence

\nopagebreak[4]
\[
J(0)=
\begin{cases}
(-1)^s4^{-s}\binom{2s}{s},&r=2s,\\[4pt]
(-1)^{s+1}2^{-2s-1}\binom{2s+1}{s},&r=2s+1.
\end{cases}
\tag{5.10}
\]

At $t=0$, Jacobi symmetry and (5.1)--(5.2) give

\nopagebreak[4]
\[
H(0)=\frac{2r}{r+2}J(0)\quad\text{if }r\text{ is even},
\]

and

\nopagebreak[4]
\[
H(0)=2J(0)\quad\text{if }r\text{ is odd}.
\]

Using (5.3), (5.6), and (5.10), we obtain

\nopagebreak[4]
\[
(-1)^rN_{[2]\times[r]\times[r+1]}(-1)
=
\begin{cases}
\displaystyle\frac{1}{s+1}\binom{2s}{s}^{\!2},
&r=2s,\\[8pt]
\displaystyle\frac{2(2s+1)}{(s+1)^2}\binom{2s}{s}^{\!2},
&r=2s+1.
\end{cases}
\tag{5.11}
\]

Both values are strictly positive, so $-1$ is not a zero. We may therefore
write

\nopagebreak[4]
\[
N_{[2]\times[r]\times[r+1]}(x)
=\prod_{i=1}^r\left((1+x)^2+\delta_i x\right),
\qquad
\delta_i=a_i+a_i^{-1}-2>0.
\tag{5.12}
\]

The coefficient $\gamma_j$ in (1.1) is now
$e_j(\delta_1,\ldots,\delta_r)$, which is strictly positive for every
$0\leq j\leq r$. This proves strict $\gamma$-positivity. Together with the
palindromicity classification, it proves Conjecture 4.5 and Conjecture B for
$k=2$, $P=[m]\times[n]$.
\end{proof}


\section{Consequences, scope, and prior-work boundary}


Real-rootedness immediately implies log-concavity and unimodality of the
coefficient sequence because all coefficients are nonnegative. Theorem 1.2 also
records the zero location explicitly: every zero lies in $(-\infty,0)$. For a
real-rooted polynomial with positive coefficients and constant term $1$, this
location is a consequence of real-rootedness rather than a logically stronger
claim. Theorem 1.3 adds a complete palindromicity classification. Its necessary
direction is already implicit in Ding and Dong's monicity classification
\cite[Lemma~2.3(a)]{ding2019}; the new part is the adjacent-parameter sufficiency and the
strict $\gamma$-positivity conclusion obtained by combining that symmetry with
Theorem 1.2 and the nonvanishing formula (5.11).

After the substitution $x=y+1$ from (2.6), Theorem 1.1 proves the determinant
of \cite{he2021}. It also explains the shifted parameters in its off-diagonal entries:
they arise from the two-unit row imbalance forced by the tail switch rather than
from numerical fitting. The proof is entirely within poset and two-rowed-array
combinatorics; the benzenoid interpretation is not needed.

The argument is specific to the two-layer factor $[2]$. For
$[k]\times[m]\times[n]$ with $k>2$, one would need a higher-order determinant or
a compatible multi-path cancellation. No such extension is claimed here. We
do not claim Conjecture B for other values of $k$ or for nonrectangular
minuscule posets.

As of 26 August 2026, targeted searches of the exact conjectures, the determinant
in \cite{he2021}, the $O(2,m,n)$ family, and the extended strict order-polynomial literature
found no independent earlier public proof of Theorems 1.1--1.3 or of the
adjacent-parameter sufficiency used in Theorem 1.3. Theorems 1.1--1.2 were
already made public in this project's v0.1 release; Theorem 1.3 is new in v0.2.
This records the search boundary, not a claim of exhaustive worldwide priority.

\section*{Acknowledgement of AI assistance}


OpenAI Codex was used as an AI-assisted tool for literature discovery, proof
stress testing, exact finite verification, and language editing. The author
reviewed the mathematical argument and assumes responsibility for the claims and
presentation.

\begin{thebibliography}{9}
\footnotesize
\raggedright

\bibitem{ding2019}
J.~Ding and C.-P.~Dong,
\emph{Antichain generating polynomials of posets},
arXiv:1905.06692v1 (2019).
\href{https://arxiv.org/abs/1905.06692}{arXiv:1905.06692}.

\bibitem{he2021}
B.-H.~He, J.~Langner, and H.~A.~Witek,
\emph{Hexagonal flakes as fused parallelograms: A determinantal formula for
Zhang--Zhang polynomials of the $O(2,m,n)$ benzenoids},
J. Chinese Chem. Soc. 68 (2021), 1231--1242.
\href{https://doi.org/10.1002/jccs.202000420}{doi:10.1002/jccs.202000420}.

\bibitem{langner2021}
J.~Langner and H.~A.~Witek,
\emph{Extended strict order polynomial of a poset and fixed elements of linear extensions},
Australas. J. Combin. 81 (2021), 187--207.
\href{https://ajc.maths.uq.edu.au/pdf/81/ajc_v81_p187.pdf}{Open-access article}.

\bibitem{langner2022}
J.~Langner and H.~A.~Witek,
\emph{ZZ polynomials of regular $m$-tier benzenoid strips as extended strict
order polynomials of associated posets---Part 1: proof of equivalence},
MATCH Commun. Math. Comput. Chem. 87 (2022), 585--620.
\href{https://doi.org/10.46793/match.87-3.585L}{doi:10.46793/match.87-3.585L}.

\bibitem{krattenthaler1995}
C.~Krattenthaler,
\emph{Counting nonintersecting lattice paths with turns},
S\'em. Lothar. Combin. 34 (1995), Article B34i, 17 pp.
\href{https://www.mat.univie.ac.at/~slc/wpapers/s34vienna.pdf}{Open-access article}.

\bibitem{rubey}
M.~Rubey,
\emph{Comment on ``Counting nonintersecting lattice paths with turns'' by
C. Krattenthaler}.
\href{https://www.mat.univie.ac.at/~slc/wpapers/s34rubey.pdf}{Open-access comment}.

\bibitem{dlmf}
NIST Digital Library of Mathematical Functions,
Chapter 18: \emph{Orthogonal Polynomials}.
\href{https://dlmf.nist.gov/18}{dlmf.nist.gov/18}.

\end{thebibliography}

\end{document}
